\documentclass[a4paper]{article}     
\usepackage{amsfonts}
\usepackage{amsmath}
\usepackage{amssymb}
\usepackage{graphicx}
\usepackage{color}

\newcommand{\Q}{\mathbb{Q}}
\newcommand{\R}{\mathbb{R}}
\newcommand{\llim}{\lim\limits}
\newcommand{\dint}{\displaystyle\int}

\begin{document}

\noindent \large Auteur: Hina Rana \\
\noindent \large Studierichting: Informatica\\
\noindent \large Oefeningenreeks 3H nr. 9.2\\

\medskip

\normalsize
$r= 2-\cos \theta$\\

Domein: $\text{dom } r = \mathbb{R}$, beperkt domein: $[0, 2\pi]$\\

Periode: $P = 2\pi$\\

Nulpunten:\\

$r = 0 \Leftrightarrow \cos\theta = 2 \Leftrightarrow$ geen oplossingen ($\emptyset$)\\

Afgeleide en extrema:\\

$r' = \sin\theta = 0 \Leftrightarrow \theta = k\pi \rightarrow \{0, \pi, 2\pi\}$\\

Verloopstabel:\\

$\begin{array}{c|ccccccc}
\theta & 0 & & \pi & & 2\pi \\
\hline
r & 1 & + & 3 & + & 1 \\
\hline
r' & 0 & + & 0 & - & 0 \\
\hline
& \text{min} & \nearrow & \text{max} & \searrow & \text{min} \\
& 1 & & 3 & & 1
\end{array}$

\begin{figure}[h]
    \centering
    \includegraphics[width=0.8\textwidth]{{3H-9.2-hina.rana.INF}.jpeg}
    \caption{Grafiek van $r = 2- \cos \theta$}
\end{figure}

\begin{thebibliography}{999}
\end{thebibliography}
\end{document} 
